\chapter{Product Measures and Change of Variables}\label{ch:products}

\begin{center}\itshape
Slice a loaf crosswise and add the areas of the slices. You get the volume. Slice it lengthwise and add again.
You get the same number. Why do the two sums agree, and what happens to the volume if the loaf is squeezed or bent?
\end{center}

The same cut does two jobs. It computes a volume, which is order. It also multiplies the laws of independent random variables, which is chance.

In 1906 Johan Ludvig Heiberg, a Danish philologist, opened a prayer book in a library in Constantinople. A thirteenth-century monk had scraped the parchment and written prayers over an older text. Under the prayers Heiberg could read Archimedes. One of the works was \emph{The Method}, which had been given up for lost. In it Archimedes explains how he found the volume of a sphere before he proved it. He cut a sphere, a cone and a cylinder into thin parallel slices, and he balanced those slices against each other on an imaginary lever.

\portrait{fubini}{Guido Fubini (1879--1943). Photograph, about 1920, author unknown. Public domain, via Wikimedia Commons.}

Guido Fubini published the modern form of that slice computation in 1907. The integral of a function on a product equals either of the two repeated integrals.

\portrait{tonelli}{Leonida Tonelli (1885--1946). Photograph, about 1920, author unknown. Public domain, via Wikimedia Commons.}

Leonida Tonelli's note of 1909 treats non-negative functions. For those, you do not add a separate integrability assumption. The two repeated integrals still agree, and both are allowed to be infinite. Both theorems ask the measures to be $\sigma$-finite. That hypothesis stays in force, and a counterexample later in the chapter shows what fails when it drops.

Bonaventura Cavalieri, in 1635, had already stated the geometric case. Solids whose slices have the same area have the same volume. His indivisibles were attacked by the Jesuits, and the dispute ran on for decades. Measure turns the principle into a theorem. The measure of a set in a product is the integral of the measures of its slices.

Earlier volumes met pieces of this story and stopped short of the general statements. Volume~1 rearranged double series of non-negative terms. Volume~4 proved the continuous rectangular integral identities it needed without importing general Tonelli or Fubini. Volume~5 named Tonelli's theorem, pointed ahead to this volume, and used the name for Ritelli's evaluation of $\zeta(2)=\frac{\pi^2}6$ and for one exchange of integrals. The same volume left the identity $\Gamma(x)\Gamma(y)=B(x,y)\Gamma(x+y)$ without a double-integral proof, because it had no change of variables. Those two calculations are Exercises~\ref{ex:basel} and~\ref{ex:betagamma}.

The second half of the chapter asks what a map does to volume. A linear map multiplies every volume by the absolute value of its determinant. That factor is what a determinant measures. A smooth bijection, with a smooth inverse, multiplies the volume of a small set near a point by the absolute value of the Jacobian there. Out of that formula come polar coordinates in every dimension, the volume of the ball of radius $1$, which tends to $0$ as the dimension grows, and the density of the Gaussian vectors from Volume~4, Chapter~6.

\section{Product \texorpdfstring{$\sigma$}{sigma}-algebras and sections}

Fix two measurable spaces $(E,\mathcal A)$ and $(F,\mathcal B)$. A \emph{measurable rectangle} is a product $A\times B$ with $A\in\mathcal A$ and $B\in\mathcal B$. The \emph{product $\sigma$-algebra} $\mathcal A\otimes\mathcal B$ is the $\sigma$-algebra on $E\times F$ generated by those rectangles. Rectangles form a $\pi$-system, because the intersection of two of them is again a rectangle:
\[
(A\times B)\cap(A'\times B')=(A\cap A')\times(B\cap B').
\]
The two projections $(x,y)\mapsto x$ and $(x,y)\mapsto y$ are measurable for $\mathcal A\otimes\mathcal B$. It is the smallest $\sigma$-algebra with that property. The product topology was the coarsest topology that made the same projections continuous (Chapter~\ref{ch:topology}). The analogy is exact, with measurable maps in place of continuous ones.

\begin{proposition}\label{prop:borelproduct}
The Borel sets of the product space are exactly the product of the Borel $\sigma$-algebras:
\[
\mathcal B(\R^m)\otimes\mathcal B(\R^n)=\mathcal B(\R^{m+n}).
\]
\end{proposition}
\begin{proof}
First take the product into the Borel sets of the big space. If $A\subset\R^m$ is open, then $A\times\R^n$ is open in $\R^{m+n}$, hence Borel. The sets $A$ for which $A\times\R^n$ is Borel form a $\sigma$-algebra, and that family contains the open sets, so it contains every Borel set. The same argument with the factors swapped shows that $\R^m\times B$ is Borel whenever $B$ is Borel. Their intersection $A\times B$ is therefore Borel. Every generator of the product $\sigma$-algebra is Borel, so
\[
\mathcal B(\R^m)\otimes\mathcal B(\R^n)\subseteq\mathcal B(\R^{m+n}).
\]
For the other inclusion, take an open set $W\subset\R^{m+n}$. Chapter~\ref{ch:topology} records that the Euclidean topology is the product topology, so $W$ is a countable union of open boxes with rational corners. Each such box is a product of an open box in $\R^m$ and an open box in $\R^n$, so it lies in the product $\sigma$-algebra. Thus every open set lies in the product, and so does every Borel set.
\end{proof}

Fix $Q\subseteq E\times F$. The \emph{section} of $Q$ at $x\in E$ is the slice
\[
Q_x=\{y\in F:(x,y)\in Q\},
\]
and the section at $y\in F$ is $Q^y=\{x\in E:(x,y)\in Q\}$. For a function $f$ on the product, the sections are the functions $f_x=f(x,\cdot)$ and $f^y=f(\cdot,y)$. Figure~\ref{fig:rw-slices} is the picture for a rectangle and for a triangle.

\begin{figure}[ht]
\centering
\fig{ch09-rw-slices}
\caption{Left: the rectangle is the product of an interval $A$ on the $x$-axis and an interval $B$ on the $y$-axis. The coral vertical slice is a copy of $B$, and the dashed horizontal slice is a copy of $A$. Right: the triangle $0\le y\le x\le 2$ is not a product. The coral slice at a fixed $x$ has length $x$, and the dashed slice at a fixed $y$ has length $2-y$. The area is the integral of either family of lengths.}
\label{fig:rw-slices}
\end{figure}

\begin{lemma}\label{lem:sections}
If $Q\in\mathcal A\otimes\mathcal B$, then $Q_x\in\mathcal B$ for every $x$ and $Q^y\in\mathcal A$ for every $y$. If $f$ is measurable for $\mathcal A\otimes\mathcal B$, then every section $f_x$ and every section $f^y$ is measurable.
\end{lemma}
\begin{proof}
Fix $x$. Let $\mathcal D$ be the family of sets $Q$ in the product whose section $Q_x$ lies in $\mathcal B$. A rectangle $A\times B$ has section $B$ when $x\in A$, and section $\varnothing$ otherwise, so every rectangle is in $\mathcal D$. Sections commute with complements and with countable unions, so $\mathcal D$ is a $\sigma$-algebra. It contains the generators, hence the whole product. The argument for $Q^y$ is the same with the factors swapped. For a function, the section of a sublevel set is the sublevel set of the section: $\{f_x\le t\}=\{f\le t\}_x$.
\end{proof}

The pattern is worth keeping. To prove a property for every measurable set, collect the sets that have the property. If that family is a $\sigma$-algebra, or a $\lambda$-system, and it contains a generating family, or a generating $\pi$-system, then it contains every measurable set. The next lemma is the same move, now with measures of slices.

\section{Product measures}

A measure $\mu$ on $(E,\mathcal A)$ is \emph{$\sigma$-finite} when $E$ is a countable union of sets of finite measure. The sets may be taken increasing. Lebesgue measure is $\sigma$-finite, every probability is $\sigma$-finite, and counting measure on a countable set is $\sigma$-finite. Counting measure on $\R$ is not: every set of finite counting measure is finite, and $\R$ is not a countable union of finite sets.

\begin{lemma}\label{lem:slicemeasurable}
Let $\mu$ and $\nu$ be $\sigma$-finite measures on $(E,\mathcal A)$ and $(F,\mathcal B)$, and let $Q\in\mathcal A\otimes\mathcal B$. Then $x\mapsto\nu(Q_x)$ is $\mathcal A$-measurable, $y\mapsto\mu(Q^y)$ is $\mathcal B$-measurable, and
\[\int_E\nu(Q_x)\,d\mu(x)=\int_F\mu(Q^y)\,d\nu(y).\]
\end{lemma}
\begin{proof}
\emph{The finite case.} Begin with the case in which both measures of the whole space are finite. Let $\mathcal D$ be the family of product-measurable $Q$ for which both slice functions are measurable and the two integrals agree. For a rectangle $Q=A\times B$, the vertical slice has measure $\nu(B)\mathbf 1_A(x)$, and the horizontal slice has measure $\mu(A)\mathbf 1_B(y)$. Both integrals equal $\mu(A)\nu(B)$. So $\mathcal D$ contains every rectangle, including $E\times F$.

If $Q\subseteq Q'$ are both in $\mathcal D$, the section of the difference is the difference of the sections. Since $\nu$ is finite, $\nu((Q'\setminus Q)_x)=\nu(Q'_x)-\nu(Q_x)$, with no subtraction of infinities (Proposition~\ref{prop:measureprops}(a)). The integrals subtract, and they are finite, so $Q'\setminus Q\in\mathcal D$. If $Q_n\in\mathcal D$ increase to $Q$, then $\nu((Q_n)_x)$ increases to $\nu(Q_x)$ by continuity from below. The limit of measurable functions is measurable, and monotone convergence passes to the limit in both integrals, so $Q\in\mathcal D$.

Thus $\mathcal D$ contains $E\times F$, it is closed under proper differences, and it is closed under increasing unions. It is a $\lambda$-system. It contains the $\pi$-system of rectangles. Dynkin's theorem (Theorem~\ref{thm:dynkin6}) puts every product-measurable set into $\mathcal D$.

\emph{The $\sigma$-finite case.} Choose $E_n\uparrow E$ and $F_n\uparrow F$ with $\mu(E_n)$ and $\nu(F_n)$ finite. The restrictions $\mu_n=\mu(\cdot\cap E_n)$ and $\nu_n=\nu(\cdot\cap F_n)$ are finite measures. The finite case, applied to $Q$, says that the truncated slice function $x\mapsto\nu_n(Q_x)\mathbf 1_{E_n}(x)$ is measurable. It is the $\nu$-measure of the section of the cut-down set $Q\cap(E_n\times F_n)$. As $n$ grows it increases to $\nu(Q_x)$, because $x$ eventually lies in $E_n$ and $\nu(Q_x\cap F_n)$ increases to $\nu(Q_x)$. The same holds with the variables swapped. The two integrals for $\mu_n$ and $\nu_n$ agree at each $n$, and monotone convergence passes to the limit.
\end{proof}

\begin{theorem}[product measure]\label{thm:product}
Assume $\mu$ and $\nu$ are $\sigma$-finite. There is a unique measure $\mu\otimes\nu$ on $\mathcal A\otimes\mathcal B$ such that $(\mu\otimes\nu)(A\times B)=\mu(A)\nu(B)$ for every measurable rectangle, with the convention $0\cdot\infty=0$. For every $Q\in\mathcal A\otimes\mathcal B$,
\[(\mu\otimes\nu)(Q)=\int_E\nu(Q_x)\,d\mu(x)=\int_F\mu(Q^y)\,d\nu(y).\]
\end{theorem}
\begin{proof}
Define $(\mu\otimes\nu)(Q)$ by the first integral. Lemma~\ref{lem:slicemeasurable} says the second integral gives the same number. The empty set has empty slices, so the measure of $\varnothing$ is $0$. If $Q$ is a countable disjoint union of sets $Q_k$, then for each $x$ the slice $Q_x$ is the disjoint union of the slices $(Q_k)_x$. Hence $\nu(Q_x)=\sum_k\nu((Q_k)_x)$. Monotone convergence lets you integrate the series term by term, so $\mu\otimes\nu$ is a measure. On a rectangle the value was computed in the lemma, and the convention $0\cdot\infty=0$ is the one already used for the integral of a simple function.

For uniqueness, let $\rho$ be any measure with the same values on rectangles. Choose $E_n\uparrow E$ and $F_n\uparrow F$ of finite measure. The rectangles $E_n\times F_n$ increase to $E\times F$ and have finite $\rho$-measure. Two measures that agree on a $\pi$-system, and that are finite on an increasing sequence from the $\pi$-system which covers the space, are equal (Theorem~\ref{thm:uniq6}).
\end{proof}

\begin{figure}[ht]
\centering
\fig{ch08-cavalieri}
\caption{Cavalieri's principle, drawn with two solids. At each height the coral slice of the straight solid has the same area as the coral slice of the slanted one. The volumes therefore agree, because each volume is the integral of those slice areas. A straight stack of coins and a slanted stack hold the same amount of metal.}
\label{fig:cavalieri}
\end{figure}

\emph{Lebesgue measure is a product.} The measures $\lambda_{m+n}$ and $\lambda_m\otimes\lambda_n$ are defined on the same $\sigma$-algebra, by Proposition~\ref{prop:borelproduct}. On a box they agree, because the volume of a box is the product of the volumes of its sides. Uniqueness in Theorem~\ref{thm:lebesgue} gives $\lambda_{m+n}=\lambda_m\otimes\lambda_n$. Iterating, $\lambda_d$ is the product of $d$ copies of $\lambda$. Theorem~\ref{thm:product} is then Cavalieri's principle, drawn in Figure~\ref{fig:cavalieri}. For a Borel set $Q\subseteq\R^d$,
\[\lambda_d(Q)=\int_\R\lambda_{d-1}(Q^t)\,dt,\qquad Q^t=\{x\in\R^{d-1}:(x,t)\in Q\}.\]

Here is the triangle in the plane, as a check you can do by hand. Let $Q$ be the set of points with $0\le x\le 1$ and $0\le y\le 1-x$, the triangle under the diagonal of the unit square. The vertical slice at $x$ has length $1-x$, so
\[
\lambda_2(Q)=\int_0^1(1-x)\,dx=\frac12.
\]
The horizontal slices give the same half. A cone is the same computation one dimension up. Let the base be a Borel set $B\subseteq\R^{d-1}$ and the height $h>0$, and let the slice at height $t\in[0,h)$ be the scaled copy $(1-t/h)B$. Scaling (Theorem~\ref{thm:lebesgue}) multiplies the $(d-1)$-dimensional measure by $(1-t/h)^{d-1}$. The volume is
\[
\lambda_{d-1}(B)\int_0^h\Bigl(1-\frac th\Bigr)^{d-1}dt=\frac hd\lambda_{d-1}(B).
\]
For $d=3$ this is the school formula: base area times height, divided by $3$. The shape of the base never entered, only its measure.

\section{Tonelli and Fubini}

Indicators are settled. A general non-negative function is a limit of simple functions, and an integrable function is a difference of two non-negative ones.

\begin{theorem}[Tonelli]\label{thm:tonelli}
Let $\mu$ and $\nu$ be $\sigma$-finite, and let $f:E\times F\to[0,+\infty]$ be $\mathcal A\otimes\mathcal B$-measurable. Then $x\mapsto\int_F f(x,y)\,d\nu(y)$ and $y\mapsto\int_E f(x,y)\,d\mu(x)$ are measurable, and
\[\int_{E\times F}f\,d(\mu\otimes\nu)=\int_E\Bigl(\int_Ff\,d\nu\Bigr)d\mu=\int_F\Bigl(\int_Ef\,d\mu\Bigr)d\nu.\]
\end{theorem}
\begin{proof}
When $f=\mathbf 1_Q$, the identity is Theorem~\ref{thm:product}. Both sides add when you add functions, and both sides scale when you multiply by a positive constant, so the identity holds for every non-negative simple function. For a general $f$, choose simple functions $0\le s_n\uparrow f$ (Proposition~\ref{prop:measurablefns}). At each fixed $x$, monotone convergence gives
\[
\int_F s_n(x,y)\,d\nu(y)\uparrow\int_F f(x,y)\,d\nu(y).
\]
The limit of measurable functions of $x$ is measurable. Apply monotone convergence once more, in the variable $x$, and once on the product space. The other order is the same argument with the factors swapped.
\end{proof}

Nothing in the statement asks $f$ to be integrable. Both sides may equal $+\infty$, and then they are infinite together. In practice you use Tonelli first, on $|f|$, to test integrability. If that integral is finite, Fubini lets you swap the order for $f$ itself.

\begin{theorem}[Fubini]\label{thm:fubini}
Let $\mu$ and $\nu$ be $\sigma$-finite, and let $f\in L^1(\mu\otimes\nu)$ be real or complex. Then for $\mu$-almost every $x$, the section $f_x$ is $\nu$-integrable. The function $x\mapsto\int_F f(x,y)\,d\nu(y)$, defined for almost every $x$, is $\mu$-integrable. The same holds with the roles of the two variables exchanged. And
\[\int_{E\times F}f\,d(\mu\otimes\nu)=\int_E\Bigl(\int_Ff\,d\nu\Bigr)d\mu=\int_F\Bigl(\int_Ef\,d\mu\Bigr)d\nu.\]
\end{theorem}
\begin{proof}
Tonelli applied to $|f|$ gives
\[
\int_E\Bigl(\int_F|f|\,d\nu\Bigr)d\mu=\int|f|\,d(\mu\otimes\nu)<\infty.
\]
The measurable function $x\mapsto\int_F|f_x|\,d\nu$ is therefore finite outside a $\mu$-null set $N$ (Proposition~\ref{prop:markov}). Suppose $f$ is real, and write $f=f^+-f^-$. Off $N$, both $\int f^+_x\,d\nu$ and $\int f^-_x\,d\nu$ are finite, so their difference is the integral of $f_x$. Each of those two functions of $x$ is measurable and integrable, by Tonelli applied to $f^+$ and to $f^-$. Integrating over $E\setminus N$ does not change the value, because $N$ is null. Tonelli on $f^+$ and on $f^-$ gives the identity. If $f$ is complex, apply the real case to the real part and to the imaginary part.
\end{proof}

\begin{think}
Can the two repeated sums of one array both converge, and converge to different numbers? Fill the diagonal of an infinite matrix with $1$, put $-1$ in the entry just to the right of each diagonal entry, and put $0$ everywhere else.
\end{think}

With counting measure on $\N$, the two theorems are statements about double series, of the kind Volume~1 treated for non-negative terms. If $a_{mn}\ge 0$, the two repeated sums agree, and the common value may be infinite. If $\sum_{m,n}|a_{mn}|<\infty$, both repeated sums of $a_{mn}$ converge, and they converge to the same number. Drop absolute convergence and the matrix in the box is a counterexample. Each row sums to $1-1=0$, so the sum of the row sums is $0$. The first column sums to $1$, and every later column sums to $-1+1=0$, so the sum of the column sums is $1$. Exercise~\ref{ex:fubinifail} is the same failure for an integral.

\begin{figure}[ht]
\centering
\fig{ch08-diagonal}
\caption{Lebesgue measure runs along the $x$-axis and counting measure runs along the $y$-axis. The diagonal $D$ of the unit square meets the coral vertical line in one point, which counting measure scores as $1$. It meets the dashed horizontal line in one point, which Lebesgue measure scores as $0$. Integrating the slices gives $1$ in one order and $0$ in the other.}
\label{fig:diagcount}
\end{figure}

\emph{$\sigma$-finiteness cannot be dropped.} Let $\mu$ be Lebesgue measure on $[0,1]$, let $\nu$ be counting measure on $[0,1]$, and let $D=\{(x,x)\}$ be the diagonal. The diagonal is closed, so Proposition~\ref{prop:borelproduct} puts it in the product of the two Borel $\sigma$-algebras. Every vertical slice is a single point, and counting measure gives it mass $1$, so $\int\nu(D_x)\,d\lambda=1$. Every horizontal slice is a single point, and Lebesgue measure gives it mass $0$, so $\int\lambda(D^y)\,d\nu=\sum_y 0=0$. Figure~\ref{fig:diagcount} is this pair of slices. The proof of Lemma~\ref{lem:slicemeasurable} used $\sigma$-finiteness exactly when it cut $\nu$ into finite pieces. Counting measure on $[0,1]$ cannot be cut that way.

\emph{Several factors, and independence.} The construction repeats. There is a unique measure $\mu_1\otimes\cdots\otimes\mu_n$ on $\mathcal A_1\otimes\cdots\otimes\mathcal A_n$ such that a rectangle of $n$ sides has measure equal to the product of the $n$ measures, and the way you insert parentheses does not matter. Both groupings agree on rectangles, so Theorem~\ref{thm:uniq6} makes them the same measure.

In probability, random variables $X_1,\dots,X_n$ are independent exactly when the law of the tuple $(X_1,\dots,X_n)$ is the product of the laws. Both candidate measures agree on rectangles, and rectangles are the sets that appear in the definition of independence. A $\sigma$-finite measure on the product is determined by its values on rectangles. Fubini then gives $\mathbb E[XY]=\mathbb E[X]\,\mathbb E[Y]$ for independent integrable $X$ and $Y$. The check that $xy$ is integrable is Tonelli:
\[
\int|xy|\,d(P_X\otimes P_Y)=\mathbb E|X|\,\mathbb E|Y|<\infty.
\]
More generally, for measurable $h\ge 0$, or for $h$ such that $h(X,Y)$ is integrable,
\[\mathbb E[h(X,Y)]=\int\mathbb E[h(x,Y)]\,dP_X(x).\]
You freeze $X$ at a value $x$, you average over $Y$, and then you average over $X$.

The \emph{convolution} $\mu*\nu$ of two probabilities on $\R^d$ is the law of $X+Y$ when $X$ and $Y$ are independent, $X$ has law $\mu$, and $Y$ has law $\nu$. Explicitly,
\[
(\mu*\nu)(A)=\iint\mathbf 1_A(x+y)\,d\mu(x)\,d\nu(y).
\]
If $\mu$ and $\nu$ have densities $f$ and $g$, a translation in the inner integral produces a density for the sum. The density is
\[(f*g)(s)=\int_{\R^d}f(s-y)\,g(y)\,dy.\]
The computation is Tonelli together with translation invariance of $\lambda_d$ (Theorem~\ref{thm:lebesgue}):
\[\begin{aligned}
\iint\mathbf 1_A(x+y)f(x)g(y)\,dx\,dy
&=\int\Bigl(\int\mathbf 1_A(s)f(s-y)\,ds\Bigr)g(y)\,dy\\
&=\int_A\Bigl(\int f(s-y)g(y)\,dy\Bigr)ds.
\end{aligned}\]
Convolution of functions, as opposed to convolution of measures, is taken up in Chapter~\ref{ch:lp}.

\section{Lebesgue measure under linear maps}

A linear map moves sets, and the question is by what factor the volume changes. Figure~\ref{fig:rw-det} is the case you should keep in mind: the unit square has area $1$, and its image has area $|\det T|$.

\begin{figure}[ht]
\centering
\fig{ch09-rw-det}
\caption{The unit square on the left has area $1$. The linear map $T(x,y)=(2x+y,\,y)$ sends it to the parallelogram spanned by the columns $(2,0)$ and $(1,1)$. Every horizontal slice of that parallelogram has length $2$, and the parallelogram runs from height $0$ to height $1$, so its area is $2$. That factor is $|\det T|$.}
\label{fig:rw-det}
\end{figure}

\begin{theorem}\label{thm:lineardet}
Let $T:\R^d\to\R^d$ be a linear map. Whenever $T$ is invertible, $\lambda_d(T(A))=|\det T|\,\lambda_d(A)$ for every Borel set $A$. Whenever $T$ is not invertible, the image $T(\R^d)$ is a null set.
\end{theorem}
\begin{proof}
Suppose $T$ is invertible. Then $T$ is a homeomorphism, so $T(A)$ is Borel whenever $A$ is. The set function $\nu(A)=\lambda_d(T(A))$ is a Borel measure. It is translation invariant, because $T(A+a)=T(A)+Ta$ and Lebesgue measure is translation invariant. Write $Q$ for the half-open unit cube $]0,1]^d$, and write $c(T)$ for the number $\nu(Q)$. Each coordinate of $Tx$ is a linear function of $x$, hence bounded on $Q$, so $T(Q)$ sits in some large box and $c(T)$ is finite. It is positive, because $T(\,]0,1[^d)$ is a nonempty open set and therefore contains a small box of positive volume. The measure $c(T)^{-1}\nu$ is translation invariant and gives $Q$ measure $1$, so Theorem~\ref{thm:lebesgue} says it equals $\lambda_d$. Thus $\nu=c(T)\lambda_d$.

If $S$ and $T$ are both invertible, then
\[
\nu_{ST}(A)=\lambda_d(S(T(A)))=c(S)\,\lambda_d(T(A))=c(S)c(T)\,\lambda_d(A),
\]
so $c(ST)=c(S)c(T)$. The determinant is multiplicative as well (Volume~1). The maps with $c(T)=|\det T|$ are therefore closed under composition. It is enough to check the equality on the three kinds of map that factor every invertible map, and then to check that every invertible map is such a product.

Those three kinds are the row operations of elimination. An invertible matrix has a nonzero entry in the first column, say in row $i$. Swap coordinates $1$ and $i$, unless the nonzero entry is already in row $1$. Scale coordinate $1$ so that the corner entry becomes $1$. For each later row, a shear clears the rest of the first column. The matrix now has first column equal to the first basis vector, and it maps the coordinate subspace $\{x_1=0\}$ to itself. The block on that subspace is an invertible map of $\R^{d-1}$. Repeat. You arrive at the identity. Reading the operations backwards, $T$ is a product of swaps, scalings of one coordinate, and shears. The inverse of each of those is a map of the same kind, so the factors of $T$ itself are of those three kinds. Both $c$ and $|\det|$ turn a product into a product, so it remains only to check the three kinds.

\emph{(i) Scale one coordinate.} The map $x_i\mapsto a x_i$ with $a\ne 0$, and the other coordinates fixed, sends the unit cube to a box with one side of length $|a|$ and the other sides of length $1$. So $c(T)=|a|$. The determinant is $a$, and $|a|=|\det T|$.

\emph{(ii) Swap two coordinates.} The unit cube is unchanged, so $c(T)=1$. A swap of two columns changes the sign of the determinant (Volume~1), so $|\det T|=1$.

\emph{(iii) A shear.} The map $x_i\mapsto x_i+a x_j$ with $i\ne j$, and the other coordinates fixed, has matrix equal to the identity plus one off-diagonal entry. That matrix is triangular once the coordinates are ordered suitably, so its determinant is $1$ (Volume~1). To see the volume, take $i=1$ and $j=2$. A point $y$ lies in the image of $]0,1]^d$ precisely when $(y_2,\dots,y_d)\in\,]0,1]^{d-1}$ and $a y_2<y_1\le a y_2+1$. Each slice in the $y_1$ direction is an interval of length $1$. Cavalieri, in the form $\lambda_d=\lambda\otimes\lambda_{d-1}$, gives
\[
c(T)=\int_{]0,1]^{d-1}}1\,d\lambda_{d-1}=1.
\]

If $T$ is not invertible, its columns are linearly dependent, so they lie in some linear hyperplane $H$. Choose a basis of the coordinate hyperplane $H_0=\{x_1=0\}$, send it to a basis of $H$, and send $e_1$ to a vector outside $H$. The resulting linear map $S$ is invertible, because its columns form a basis of $\R^d$, and $S(H_0)=H$. The hyperplane $H_0$ is the product $\{0\}\times\R^{d-1}$. Theorem~\ref{thm:product}, with the convention $0\cdot\infty=0$, gives
\[
\lambda_d(H_0)=\lambda(\{0\})\,\lambda_{d-1}(\R^{d-1})=0.
\]
Hence $\lambda_d(H)=|\det S|\cdot 0=0$. The image of $T$ is contained in $H$, so it is null.
\end{proof}

Rotations and reflections have determinant $\pm 1$, so they preserve Lebesgue measure. That is the fact Chapter~\ref{ch:measures} postponed. The parallelepiped spanned by vectors $v_1,\dots,v_d$ is the image of the unit cube under the linear map with those columns, so its volume is $|\det(v_1,\dots,v_d)|$. Up to sign, the determinant is the factor by which the linear map multiplies volumes.

\section{The change of variables formula}

Linear maps multiply every set by the same factor. A curved map does not. On a small cube it is close to its derivative, and the derivative multiplies volumes by the absolute value of the Jacobian. Adding those contributions is the change of variables formula. Figure~\ref{fig:diffeo} is one square and the parallelogram that approximates its image.

A map $\varphi$ from an open set $U\subseteq\R^d$ into $\R^d$ is \emph{$C^1$} when each coordinate function $\varphi_i$ has partial derivatives, and those partial derivatives are continuous. The \emph{Jacobian matrix} $D\varphi(x)$ has entry $(i,j)$ equal to $\partial_j\varphi_i(x)$, the ordinary derivative of $t\mapsto\varphi_i(x+te_j)$ at $t=0$. The \emph{Jacobian} is the determinant $J_\varphi(x)=\det D\varphi(x)$.

One chain rule is needed, and it is only the chain rule along a curve. Let $\psi$ be $C^1$ on an open set, and let $\gamma$ be a $C^1$ curve whose image lies in that open set. The claim is
\[
(\psi\circ\gamma)'(t)=D\psi(\gamma(t))\,\gamma'(t).
\]
To see it, write the change in the $k$-th component of $\psi$ from $\gamma(t)$ to $\gamma(t+s)$ as a sum of $d$ increments, changing one coordinate at a time. On the $i$-th increment the mean value theorem (Theorem~\ref{thm:mvt}) supplies a point, on the segment just travelled, at which the difference equals $\partial_i\psi_k$ times the change in that coordinate. For small $s$ the segments stay in the domain, because the domain is open. Divide by $s$. As $s\to 0$, the intermediate point tends to $\gamma(t)$, the partial derivatives are continuous, and the difference quotient of $\gamma_i$ tends to $\gamma_i'(t)$. That is the $k$-th row of the matrix product.

Now let $\varphi:U\to V$ be a bijection of open sets with $\varphi$ and $\varphi^{-1}$ both $C^1$. Such a map is a \emph{$C^1$ diffeomorphism}. Apply the curve rule to $\psi=\varphi^{-1}$ and to the curve $t\mapsto\varphi(x+te_j)$. The composition is the straight line $t\mapsto x+te_j$, whose derivative is the basis vector $e_j$. The same curve rule, applied to $\varphi$ along that straight line, says that the velocity of $t\mapsto\varphi(x+te_j)$ at $t=0$ is the $j$-th column of $D\varphi(x)$. The right side of the rule for $\varphi^{-1}$ is therefore the $j$-th column of the product $D\varphi^{-1}(\varphi(x))\,D\varphi(x)$. That column equals $e_j$. So
\[
D\varphi^{-1}(\varphi(x))\,D\varphi(x)=I.
\]
Taking determinants, $J_\varphi(x)\,J_{\varphi^{-1}}(\varphi(x))=1$. In particular $J_\varphi(x)\ne 0$. Near $x$, the map $h\mapsto\varphi(x+h)$ is close to the affine map $h\mapsto\varphi(x)+D\varphi(x)h$, and that affine map multiplies volumes by $|J_\varphi(x)|$ (Theorem~\ref{thm:lineardet}).

\begin{figure}[ht]
\centering
\fig{ch08-diffeo}
\caption{A grid of small squares, and the grid after a diffeomorphism $\varphi$. The coral square is carried to a curved cell. The dashed parallelogram is the one spanned by the columns of $D\varphi(x)$, scaled by the side of the square. Its area is $|J_\varphi(x)|$ times the area of the square, which is why the Jacobian appears in the change of variables.}
\label{fig:diffeo}
\end{figure}

\begin{theorem}[change of variables]\label{thm:cov}
Let $\varphi:U\to V$ be a $C^1$ diffeomorphism between open subsets of $\R^d$. For every Borel function $f:V\to[0,+\infty]$,
\[\int_Vf(y)\,dy=\int_Uf(\varphi(x))\,|J_\varphi(x)|\,dx.\]
The same identity holds for every $f\in L^1(V)$, and then $(f\circ\varphi)\,|J_\varphi|$ belongs to $L^1(U)$. In particular, $\lambda_d(\varphi(A))=\int_A|J_\varphi|\,d\lambda_d$ for every Borel set $A\subseteq U$.
\end{theorem}

The arguments below measure distance by the largest absolute coordinate, $\|x\|=\max_i|x_i|$. Closed balls for this norm are cubes. The cube of centre $a$ and side $2h$ is the set of $x$ with $\|x-a\|\le h$, and its volume is $(2h)^d$. Two points of that cube are at most $2h$ apart. A matrix $M$ gets the norm equal to its largest absolute row sum, $\|M\|=\max_i\sum_j|M_{ij}|$. Then $\|Mx\|\le\|M\|\,\|x\|$, the norm multiplies under matrix products, and $\|I\|=1$.

\begin{lemma}\label{lem:cube}
Let $Q\subseteq U$ be a closed cube of centre $a$ and side $2h$, and let $T$ be an invertible matrix. The measure of the image is at most $|\det T|$ times a power of the largest stretch of $T^{-1}D\varphi$ on the cube, times the volume of the cube:
\[\lambda_d(\varphi(Q))\le|\det T|\Bigl(\sup_{z\in Q}\|T^{-1}D\varphi(z)\|\Bigr)^d\lambda_d(Q).\]
\end{lemma}
\begin{proof}
Replace $\varphi$ by the adjusted map $g=T^{-1}\circ\varphi$, so that the linear factor $T$ can be put back at the end by Theorem~\ref{thm:lineardet}. Each component of $g$ is a fixed linear combination of the components of $\varphi$, so $Dg(z)=T^{-1}D\varphi(z)$. Let $M=\sup_{z\in Q}\|Dg(z)\|$. The supremum is finite because $Dg$ is continuous on the compact set $Q$. The cube is convex, so the segment from $a$ to a point $x\in Q$ stays in $Q$. Apply the mean value theorem (Theorem~\ref{thm:mvt}) to $t\mapsto g_i(a+t(x-a))$ on $[0,1]$. For some $z$ on that segment,
\[
|g_i(x)-g_i(a)|=\Bigl|\sum_j\partial_j g_i(z)\,(x_j-a_j)\Bigr|\le M\|x-a\|\le Mh.
\]
Thus $g(Q)$ sits in the cube of centre $g(a)$ and side $2Mh$, whose volume is $M^d\lambda_d(Q)$. Theorem~\ref{thm:lineardet} gives
\[
\lambda_d(\varphi(Q))=\lambda_d(T(g(Q)))=|\det T|\,\lambda_d(g(Q))\le|\det T|\,M^d\lambda_d(Q).
\]
\end{proof}

\begin{proof}[Proof of Theorem~\ref{thm:cov}]
\emph{Step 1, on closed cubes.} Take a closed cube $Q$ contained in $U$. The aim of this step is an upper bound: the measure of the image is at most the integral of $|J_\varphi|$ over the cube. On the compact set $Q$ the matrix $D\varphi$ is continuous, and $J_\varphi$ does not vanish, so the inverse matrix $D\varphi(x)^{-1}$ exists at every point. Its entries are quotients of determinants by $J_\varphi(x)$ (Volume~1), hence continuous, and therefore bounded: $\|D\varphi(x)^{-1}\|\le C$ for some finite $C$. Heine's theorem (Theorem~\ref{thm:heine}) gives uniform continuity of $x\mapsto D\varphi(x)$ on $Q$.

Fix $\varepsilon>0$, and choose $\delta>0$ so that $\|D\varphi(z)-D\varphi(x)\|\le\varepsilon$ whenever $x,z\in Q$ and $\|z-x\|\le\delta$. Cut $Q$ into finitely many closed subcubes $Q_k$, of centres $x_k$, whose diameters in the norm $\|\cdot\|$ are at most $\delta$. Neighbouring subcubes meet on faces. A face lies in a coordinate hyperplane, and those are null (Chapter~\ref{ch:measures}), so the sum of the volumes $\lambda_d(Q_k)$ equals $\lambda_d(Q)$. For $z\in Q_k$,
\begin{align*}
\|D\varphi(x_k)^{-1}D\varphi(z)\|
&\le\|I\|+\|D\varphi(x_k)^{-1}\|\,\|D\varphi(z)-D\varphi(x_k)\|\\
&\le 1+C\varepsilon.
\end{align*}
Lemma~\ref{lem:cube}, with $T=D\varphi(x_k)$, and countable subadditivity give
\[
\lambda_d(\varphi(Q))\le\sum_k\lambda_d(\varphi(Q_k))\le(1+C\varepsilon)^d\sum_k|J_\varphi(x_k)|\,\lambda_d(Q_k).
\]
The function $|J_\varphi|$ is continuous on $Q$, hence uniformly continuous. Let $\omega(\delta)$ be the maximum of $\bigl||J_\varphi(z)|-|J_\varphi(w)|\bigr|$ over pairs in $Q$ at distance at most $\delta$. On $Q_k$ the integrand $|J_\varphi|$ stays within $\omega(\delta)$ of the constant $|J_\varphi(x_k)|$, and the faces do not affect the integral, so
\[
\Bigl|\sum_k|J_\varphi(x_k)|\,\lambda_d(Q_k)-\int_Q|J_\varphi|\Bigr|\le\omega(\delta)\,\lambda_d(Q).
\]
Hold $\varepsilon$ fixed and let the mesh $\delta$ tend to $0$. The oscillation $\omega(\delta)$ tends to $0$ with it, and the image measure is at most $(1+C\varepsilon)^d$ times the integral of $|J_\varphi|$ over $Q$. Finally let $\varepsilon$ tend to $0$. The factor in front tends to $1$, which is the bound claimed for the cube.

\emph{Step 2, on Borel sets.} Let $W\subseteq U$ be open. Among the closed dyadic cubes of side at most $1$ that are contained in $W$, keep those that are maximal under inclusion. Two dyadic cubes are nested or else have disjoint interiors, so the maximal ones have disjoint interiors. Every point of $W$ lies in some dyadic cube of small side contained in $W$, and that cube sits inside a maximal one, so the maximal cubes $Q_k$ cover $W$. Their boundaries are null, as in Step~1. Step~1 and subadditivity give
\[
\lambda_d(\varphi(W))\le\sum_k\lambda_d(\varphi(Q_k))\le\sum_k\int_{Q_k}|J_\varphi|=\int_W|J_\varphi|.
\]
Write $U$ as a union of open sets $U_n$, each bounded, with compact closure still inside $U$. One choice is
\[
U_n=\Bigl\{x\in U:\|x\|<n,\ \mathrm{dist}(x,\R^d\setminus U)>\frac1n\Bigr\}.
\]
On each $U_n$ the continuous function $|J_\varphi|$ is bounded. Now take a Borel set $A$ contained in $U_n$. Outer regularity (Proposition~\ref{prop:regularity}) gives open sets $W_m$ with $A\subseteq W_m\subseteq U_n$ and $\lambda_d(W_m\setminus A)\to 0$. Replacing $W_m$ by the intersection of the first $m$ of them, we may take the sequence decreasing. The image of $A$ is smaller than the image of $W_m$, and Step~2 bounds the latter by the integral of $|J_\varphi|$ over $W_m$. The integrands $|J_\varphi|\mathbf 1_{W_m}$ decrease to $|J_\varphi|\mathbf 1_A$ almost everywhere, because the intersection of the $W_m$ exceeds $A$ by a null set. They are dominated by $|J_\varphi|\mathbf 1_{U_n}$, which is integrable: on $U_n$ the Jacobian is bounded and the set has finite measure. Dominated convergence (Theorem~\ref{thm:dct6}) gives $\int_{W_m}|J_\varphi|\to\int_A|J_\varphi|$.

For a general Borel set $A\subseteq U$, split $A$ into the disjoint pieces $A\cap(U_n\setminus U_{n-1})$, with $U_0=\varnothing$, and add the inequalities. Adding them shows that the measure of the image is at most the integral of $|J_\varphi|$ over $A$.

\emph{Step 3, from sets to functions.} Let $B\subseteq V$ be Borel. Step~2 applied to $A=\varphi^{-1}(B)$ reads
\[
\int_V\mathbf 1_B=\lambda_d(B)\le\int_U\mathbf 1_B(\varphi(x))\,|J_\varphi(x)|\,dx.
\]
Linear combinations give the same inequality for every non-negative simple function. Monotone convergence extends it to every Borel $f\ge 0$:
\[
\int_V f\le\int_U(f\circ\varphi)\,|J_\varphi|.
\]
Call this the forward estimate.

\emph{Step 4, running the estimate backwards.} Apply the forward estimate to the diffeomorphism $\varphi^{-1}:V\to U$ and to the non-negative function $g=(f\circ\varphi)\,|J_\varphi|$ on $U$. You get
\[
\int_U g\le\int_V(g\circ\varphi^{-1})\,|J_{\varphi^{-1}}|.
\]
The integrand on the right is
\[
f(y)\,|J_\varphi(\varphi^{-1}(y))|\,|J_{\varphi^{-1}}(y)|.
\]
The product of the two derivative matrices is the identity, so the product of the Jacobians equals $1$, and the product of their absolute values equals $1$ as well. So $\int_U g\le\int_V f$. Together with the forward estimate, this is the identity for $f\ge 0$.

If $f$ is integrable, the identity for $|f|$ shows that $(f\circ\varphi)\,|J_\varphi|$ is integrable. Apply the non-negative case to the positive and negative parts of the real and imaginary parts.
\end{proof}

Two remarks follow the identity, and neither needs a new argument. First, $\varphi$ has to be bijective only up to null sets. If $U\setminus U'$ and $V\setminus V'$ are null and $\varphi:U'\to V'$ is a diffeomorphism, the integrals over $U$ and over $U'$ agree, and the integrals over $V$ and over $V'$ agree. Second, in dimension one, for $\varphi$ increasing, the identity is the substitution rule of Volume~1, now valid for every non-negative or integrable Borel function.

\section{Polar coordinates, the \texorpdfstring{$n$}{n}-ball and Gaussian vectors}

\emph{The plane, first.} The map $(r,\theta)\mapsto(r\cos\theta,\,r\sin\theta)$ is a diffeomorphism from $]0,\infty[\times\,]0,2\pi[$ onto the plane with the half-line $[0,\infty[\times\{0\}$ removed. That half-line is null. The Jacobian is
\[
\cos\theta\cdot r\cos\theta-(-r\sin\theta)\sin\theta=r.
\]
Theorem~\ref{thm:cov} therefore gives, for $f\ge 0$,
\[
\int_{\R^2}f=\int_0^{2\pi}\int_0^\infty f(r\cos\theta,\,r\sin\theta)\,r\,dr\,d\theta.
\]
Take $f(x,y)=e^{-x^2-y^2}$. Tonelli separates the square of the one-dimensional integral:
\[\Bigl(\int_\R e^{-x^2}dx\Bigr)^2=\int_0^{2\pi}\!\!\int_0^\infty e^{-r^2}r\,dr\,d\theta=2\pi\cdot\frac12=\pi,\]
so $\int_\R e^{-x^2}\,dx=\sqrt\pi$. This is Poisson's evaluation. The earlier volumes used it. It is a theorem now.

In higher dimension one can write spherical coordinates by hand. Exercise~\ref{ex:spherical3} does that in $\R^3$. A cleaner split keeps the radius and the direction apart. Let $S=S^{d-1}=\{x\in\R^d:|x|=1\}$, with the Euclidean norm, and set
\[
\Phi(x)=\Bigl(|x|,\,\frac{x}{|x|}\Bigr).
\]
This is a homeomorphism from $\R^d\setminus\{0\}$ onto $]0,\infty[\times S$, with inverse $(r,\omega)\mapsto r\omega$.

\begin{theorem}[polar coordinates]\label{thm:polar}
There is a unique Borel measure $\sigma$ on $S^{d-1}$ such that for every Borel function $f:\R^d\to[0,+\infty]$,
\[\int_{\R^d}f(x)\,dx=\int_0^\infty\!\!\int_{S^{d-1}}f(r\omega)\,d\sigma(\omega)\,r^{d-1}\,dr.\]
It is given by $\sigma(B)=d\cdot\lambda_d(\{r\omega:0<r\le 1,\ \omega\in B\})$, and it is invariant under rotations.
\end{theorem}
\begin{proof}
Define $\sigma$ by that formula. If the sets $B_k$ are disjoint, the corresponding cones are disjoint, so $\sigma$ is a measure. It is finite because the cone over the whole sphere, out to radius $1$, sits inside the unit ball. Let $\rho$ be the measure on the half-line with density $r^{d-1}$ against one-dimensional Lebesgue measure. Compare the image measure $\Phi_*\lambda_d$, that is the set function $A\mapsto\lambda_d(\Phi^{-1}(A))$, with the product $\rho\otimes\sigma$.

On a rectangle $]a,b]\times B$ with $0\le a<b$, the preimage is the cone of $B$ out to radius $b$, with the cone out to radius $a$ removed. Scaling (Theorem~\ref{thm:lebesgue}) multiplies volumes by the $d$-th power of the factor, so
\[
\lambda_d\bigl(\Phi^{-1}(]a,b]\times B)\bigr)=(b^d-a^d)\frac{\sigma(B)}d.
\]
The integral of $r^{d-1}$ from $a$ to $b$ equals $(b^d-a^d)/d$, which is $\rho(]a,b])$. The two measures therefore agree on these rectangles.

The rectangles generate the Borel $\sigma$-algebra of $]0,\infty[\times S$. To see that, build countable bases. The half-line has the open intervals with rational endpoints. The sphere is closed and bounded in $\R^d$, hence compact (Theorem~\ref{thm:heineborel}). For each $n$, finitely many balls of radius $1/n$ cover it. Let $D$ be the set of all those centres, as $n$ runs. The balls of rational radius centred at points of $D$ form a countable collection, and every open subset of $S$ is a union of members of that collection. Products of one basis element from each factor form a countable basis of the product topology. Every open set in the product is a countable union of such products, hence lies in the product $\sigma$-algebra. So every Borel set of the product lies in $\rho$'s and $\sigma$'s product $\sigma$-algebra. In the other direction, the coordinate projections are continuous (Chapter~\ref{ch:topology}), so the preimage of a Borel set under either projection is Borel in the product. The two $\sigma$-algebras agree.

Both measures are finite on each rectangle $]\frac1n,n]\times S$, and these rectangles increase to the whole half-line times the sphere. Theorem~\ref{thm:uniq6} identifies the image of Lebesgue measure with the product $\rho\otimes\sigma$ on that space. The origin is a single point, hence null, so
\[
\int_{\R^d}f\,d\lambda_d=\int(f\circ\Phi^{-1})\,d(\rho\otimes\sigma).
\]
Tonelli is the stated formula.

For uniqueness, put $f=\mathbf 1_{\Phi^{-1}(]0,1]\times B)}$ into the formula. The right side equals $\sigma(B)/d$, and the left side is the volume of the cone, so $\sigma$ is forced. A rotation $R$ preserves $\lambda_d$, because $|\det R|=1$ (Theorem~\ref{thm:lineardet}), and it sends the cone of $B$ to the cone of $RB$. Hence $\sigma(RB)=\sigma(B)$.
\end{proof}

A radial function depends on a point only through its Euclidean norm. The formula collapses to one integral:
\[
\int_{\R^d}g(|x|)\,dx=\sigma(S^{d-1})\int_0^\infty g(r)\,r^{d-1}\,dr.
\]
Apply this to $g(r)=e^{-r^2}$. Tonelli and the one-dimensional Gaussian integral give $\int e^{-|x|^2}\,dx=(\sqrt\pi)^d$. The radial integral is a Gamma integral after the substitution $t=r^2$. The change of variables in one dimension replaces $r^{d-1}\,dr$ by $\frac12 t^{d/2-1}\,dt$, so
\[
\int_0^\infty e^{-r^2}r^{d-1}\,dr=\frac12\int_0^\infty t^{d/2-1}e^{-t}\,dt=\frac12\Gamma\Bigl(\frac d2\Bigr),
\]
where $\Gamma(s)=\int_0^\infty t^{s-1}e^{-t}\,dt$. Therefore
\[\pi^{d/2}=\sigma(S^{d-1})\int_0^\infty e^{-r^2}r^{d-1}\,dr=\sigma(S^{d-1})\,\frac{\Gamma(\frac d2)}2.\]

\begin{theorem}[the $n$-ball]\label{thm:ball}
The unit sphere and the unit ball of $\R^d$ have measures
\[\sigma(S^{d-1})=\frac{2\pi^{d/2}}{\Gamma(\frac d2)},\qquad\lambda_d(B^d)=\frac{\sigma(S^{d-1})}d=\frac{\pi^{d/2}}{\Gamma(\frac d2+1)}.\]
\end{theorem}
\begin{proof}
The formula for $\sigma(S^{d-1})$ is the identity just obtained, solved for the surface measure. For the ball, the radial formula with $g=\mathbf 1_{[0,1]}$ gives
\[
\lambda_d(B^d)=\sigma(S^{d-1})\int_0^1 r^{d-1}\,dr=\frac{\sigma(S^{d-1})}d.
\]
The Gamma relation $z\Gamma(z)=\Gamma(z+1)$, at $z=d/2$, rewrites the quotient as $\pi^{d/2}/\Gamma(d/2+1)$.
\end{proof}

The volumes in dimensions $1$, $2$, $3$ and $4$ are $2$, $\pi$, $\frac43\pi$ and $\frac{\pi^2}2$. You can check each one from the formula. In dimension $1$ the unit ball is the interval $[-1,1]$, of length $2$. In dimension $3$ the volume is the familiar $\frac43\pi$. This is the formula of Volume~5, Chapter~5. There the slicing recursion was taken as the definition of the volume. Here that recursion is Cavalieri's principle applied to $\lambda_d$, so the number computed in Volume~5 is the Lebesgue measure of the ball. Since $\sigma(S^1)=2\pi$ is the length of the circle and $\sigma(S^2)=4\pi$ is the area of the sphere, the measure $\sigma$ is surface measure.

\begin{think}
The unit ball in dimension $100$ has a volume given by the formula above. How small is it beside the cube $[-1,1]^{100}$ that contains it? Where in that cube does the volume sit?
\end{think}

The formula gives
\[
\lambda_{100}(B^{100})=\frac{\pi^{50}}{50!}\approx 2.4\cdot 10^{-40},
\]
while the cube has volume $2^{100}\approx 1.3\cdot 10^{30}$. The ball occupies a fraction of about $2\cdot 10^{-70}$ of the cube. Almost all of the cube's volume sits in the corners. The average of $x_1^2+\cdots+x_{100}^2$, for $x$ uniform in the cube, equals $100/3$: each coordinate is uniform on $[-1,1]$, and
\[
\frac12\int_{-1}^1 x^2\,dx=\frac13.
\]
The square root of $100/3$ is about $5.8$, so a typical point of the cube is that far from the centre, far outside the unit ball. Inside the ball the volume hugs the surface. The ball of radius $0.9$ has fraction $0.9^{100}\approx 2.7\cdot 10^{-5}$ of the volume of the unit ball. High-dimensional geometry does not look like geometry in dimension three. This mismatch is called the curse of dimensionality. If you sample the cube uniformly and keep a point only when it lands in the ball, you almost never keep one. An average that samples the coordinates one at a time does not need that hit.

\emph{Gaussian vectors.} A standard Gaussian vector $Z=(Z_1,\dots,Z_d)$ has independent coordinates, each of density $e^{-t^2/2}/\sqrt{2\pi}$. The joint density is the product
\[
\prod_{i=1}^d\frac{e^{-z_i^2/2}}{\sqrt{2\pi}}=(2\pi)^{-d/2}e^{-|z|^2/2},
\]
which depends on $z$ only through $|z|$. A \emph{Gaussian vector} is a random vector $X=AZ+m$ with $A$ a $d\times d$ matrix and $m\in\R^d$. Its covariance matrix is $\Sigma=AA^{\mathsf T}$.

\begin{theorem}[Gaussian densities]\label{thm:gaussdensity}
If $\Sigma=AA^{\mathsf T}$ is invertible, then $X=AZ+m$ has density
\[f_X(x)=\frac1{(2\pi)^{d/2}\sqrt{\det\Sigma}}\exp\Bigl(-\frac12(x-m)^{\mathsf T}\Sigma^{-1}(x-m)\Bigr).\]
In particular the law of $X$ depends only on $m$ and $\Sigma$. If $R$ is orthogonal, then $RZ$ is again standard Gaussian.
\end{theorem}
\begin{proof}
The identity $\det\Sigma=(\det A)^2$ is nonzero, so $A$ is invertible. The map $\varphi(z)=Az+m$ is a diffeomorphism of $\R^d$, and $|J_\varphi|=|\det A|=\sqrt{\det\Sigma}$. For a Borel set $B$, Theorem~\ref{thm:cov} applied to $\varphi^{-1}$ gives
\begin{align*}
\mathbb P(X\in B)
&=\int\mathbf 1_B(Az+m)\,\frac{e^{-|z|^2/2}}{(2\pi)^{d/2}}\,dz\\
&=\int_B\frac{e^{-|A^{-1}(x-m)|^2/2}}{(2\pi)^{d/2}\,|\det A|}\,dx.
\end{align*}
The quadratic form in the exponent is
\[
|A^{-1}(x-m)|^2=(x-m)^{\mathsf T}(A^{-1})^{\mathsf T}A^{-1}(x-m)=(x-m)^{\mathsf T}\Sigma^{-1}(x-m).
\]
The density is the integrand. It mentions $A$ only through $m$ and $\Sigma$. If $R$ is orthogonal, then $RR^{\mathsf T}=I$, so $RZ$ has mean $0$ and covariance $I$, and the density just proved is the standard Gaussian density.
\end{proof}

Volume~4, Chapter~6, described Gaussian vectors by their characteristic functions. The formula above is their density. The radius $|Z|$ and the direction $Z/|Z|$ split apart in polar coordinates, which is why the direction is uniform on the sphere (Exercise~\ref{ex:sphere}). The same split, in the plane, is what Box and Muller used in 1958 to build a pair of independent standard normals from two uniform numbers.

\begin{beauty}
One identity, $\int_V f=\int_U(f\circ\varphi)\,|J_\varphi|$, contains the substitution rule of Volume~1, polar coordinates, the volume of the ball in every dimension, the Gaussian integral, the density of every non-degenerate Gaussian vector, and the geometric meaning of the determinant. Two facts sit under it. Lebesgue measure is the only translation-invariant measure on Borel sets that gives the unit cube measure $1$. On a small enough cube, a $C^1$ map is close to a linear map.
\end{beauty}

\section{What we proved, and what comes next}

Two measures have a product (Theorem~\ref{thm:product}), and the integral for the product is an iterated integral, in either order: always for non-negative
functions (Tonelli, Theorem~\ref{thm:tonelli}), and for integrable ones (Fubini, Theorem~\ref{thm:fubini}). This pays the debt that Volumes~4 and~5 recorded under the name of Tonelli.
A linear map multiplies Lebesgue measure by the absolute value of its determinant (Theorem~\ref{thm:lineardet}), a diffeomorphism by its Jacobian
(Theorem~\ref{thm:cov}). From these came polar coordinates in every dimension (Theorem~\ref{thm:polar}), the volume of the ball (Theorem~\ref{thm:ball}) and the density of
a Gaussian vector (Theorem~\ref{thm:gaussdensity}). Nothing in the chapter was admitted. The two counterexamples, the diagonal with counting measure and
Exercise~\ref{ex:fubinifail}, show that neither $\sigma$-finiteness nor integrability can be dropped.

So far a function has been a thing to integrate. In the next chapter functions become points. The integrable functions form a space with a
distance, the square-integrable ones a space with an inner product, and the first question about such a space is the one Chapter~\ref{ch:metric} asked
about the rationals: is it complete? Chapter~\ref{ch:lp} shows that it is, and that is the reason Lebesgue's integral replaced Riemann's.

\section{Exercises}

\begin{exercise}[$\star$]\label{ex:fubinifail}
On the open square $]0,1[^2$, let $f(x,y)=\frac{x^2-y^2}{(x^2+y^2)^2}$. Show that
\[
\int_0^1\Bigl(\int_0^1 f(x,y)\,dy\Bigr)dx=\frac\pi4
\]
and that the integral in the other order equals $-\frac\pi4$. Which hypothesis of Fubini's theorem is the one that fails? Use polar coordinates to see the failure directly.
\end{exercise}

\begin{exercise}[$\star\star$]\label{ex:basel}
\emph{Ritelli's evaluation of $\zeta(2)=\frac{\pi^2}6$, from Volume~5, Chapter~3.} Set
\[
I=\int_0^\infty\!\int_0^\infty\frac{dx\,dy}{(1+y)(1+x^2y)}
\]
and compute $I$ in both orders. Deduce $\int_0^1\frac{-\ln x}{1-x^2}\,dx=\frac{\pi^2}8$, then $\sum_{k\ge 0}\frac1{(2k+1)^2}=\frac{\pi^2}8$, and finally $\zeta(2)=\frac{\pi^2}6$. Every time you swap a sum and an integral, say why the swap is allowed.
\end{exercise}

\begin{exercise}[$\star\star$]\label{ex:betagamma}
For $a>0$ and $b>0$, the Beta integral is the integral of $v^{a-1}(1-v)^{b-1}$ from $0$ to $1$. Call that number $B(a,b)$. Prove that the product of the Gamma values equals this Beta value times $\Gamma(a+b)$. Start from the product of the two Gamma integrals over $]0,\infty[^2$, and change variables by $(s,t)\mapsto(u,v)=\bigl(s+t,\frac{s}{s+t}\bigr)$.
\end{exercise}

\begin{exercise}[$\star$]\label{ex:gammasum}
The Gamma law of parameter $a>0$ is the probability on the positive half-line with density $t^{a-1}e^{-t}/\Gamma(a)$. Show that the sum of two independent Gamma random variables, of parameters $a$ and $b$, is Gamma of parameter $a+b$. Deduce the law of the sum of $n$ independent exponential random variables of parameter $1$.
\end{exercise}

\begin{exercise}[$\star\star$]\label{ex:sphere}
\emph{Uniform measure on the sphere.} Let $Z$ be a standard Gaussian vector in $\R^d$. Show that the radius $|Z|$ and the direction $Z/|Z|$ are independent, and that the direction has law $\sigma/\sigma(S^{d-1})$. If $U$ is uniform on $]0,1[$ and independent of $Z$, show that $U^{1/d}Z/|Z|$ is uniform in the unit ball.
\end{exercise}

\begin{exercise}[$\star$]\label{ex:spherical3}
Show that
\[
(r,\theta,\phi)\mapsto(r\sin\theta\cos\phi,\,r\sin\theta\sin\phi,\,r\cos\theta)
\]
is a diffeomorphism from $]0,\infty[\times\,]0,\pi[\times\,]0,2\pi[$ onto $\R^3$ minus a null set, with Jacobian $r^2\sin\theta$. Recover the volume of the unit ball and the area of the unit sphere, and compute $\int_{\R^3}\frac{e^{-|x|}}{|x|}\,dx$.
\end{exercise}

\begin{exercise}[$\star\star$]\label{ex:shell}
\emph{Where the volume sits.} (a) Fix $\varepsilon>0$. Show that the fraction of the volume of $B^d$ lying at distance less than $\varepsilon$ from the sphere tends to $1$ as $d\to\infty$. (b) Show that the fraction of the volume of $B^d$ lying in the slab $|x_1|\le c/\sqrt d$ tends to
\[
\int_{-c}^{c}e^{-t^2/2}\,\frac{dt}{\sqrt{2\pi}}.
\]
So in high dimension, almost all the volume of the ball is near the sphere, and near every equator.
\end{exercise}

\begin{exercise}[$\star\star\star$]\label{ex:simplex}
\emph{Dirichlet's integral.} Let $\Delta_n$ be the simplex of positive $n$-tuples whose coordinates sum to less than $1$, and let $a_1,\dots,a_{n+1}>0$. Show that
\[\int_{\Delta_n}x_1^{a_1-1}\cdots x_n^{a_n-1}(1-x_1-\cdots-x_n)^{a_{n+1}-1}\,dx=\frac{\Gamma(a_1)\cdots\Gamma(a_{n+1})}{\Gamma(a_1+\cdots+a_{n+1})},\]
by evaluating the product of the Gamma integrals over the positive orthant in the coordinates $s=t_1+\cdots+t_{n+1}$ and $x_i=t_i/s$ for $i\le n$. Deduce that the simplex has $n$-dimensional volume $1/n!$.
\end{exercise}

\begin{exercise}[$\star$]\label{ex:frullani}
\emph{Two short computations.} (a) For $a,b,c>0$, find the volume of the ellipsoid $\frac{x^2}{a^2}+\frac{y^2}{b^2}+\frac{z^2}{c^2}\le 1$ without evaluating an integral. (b) Let $0<a<b$. Integrate $e^{-xy}$ over $]0,\infty[\times[a,b]$ in both orders, and deduce
\[
\int_0^\infty\frac{e^{-ax}-e^{-bx}}{x}\,dx=\ln\frac ba.
\]
\end{exercise}

\begin{exercise}[$\star$]\label{ex:v7c09a}
Let $Q=\{(x,y)\in\,]0,1[^2:y<x^2\}$. Compute $\int_0^1\bigl(\int_0^1\mathbf 1_Q(x,y)\,dy\bigr)dx$ and the repeated integral in the other order. Which theorem of this chapter says the two numbers agree before you calculate them?
\end{exercise}

\begin{exercise}[$\star\star$]\label{ex:v7c09b}
Let $T:\R^2\to\R^2$ be the linear map $T(x,y)=(2x+y,\,y)$, as in Figure~\ref{fig:rw-det}, and let $S=[0,1]^2$. Show that $T(S)$ is the set of points whose second coordinate $v$ lies in $[0,1]$ and whose first coordinate lies between $v$ and $v+2$, and that this set has area $2$. Get the area once from Theorem~\ref{thm:lineardet}, and once by integrating the lengths of the horizontal slices.
\end{exercise}

\begin{exercise}[$\star\star$]\label{ex:v7c09c}
Let $U$ be the open rectangle of radii from $1$ to $2$ and angles from $0$ to $\pi/2$, and let $\varphi(r,\theta)=(r\cos\theta,\,r\sin\theta)$. Show that $\varphi$ is a $C^1$ diffeomorphism of $U$ onto the open quarter annulus
\[
\{(x,y):x>0,\ y>0,\ 1<x^2+y^2<4\},
\]
and that $J_\varphi(r,\theta)=r$. Deduce that this quarter annulus has area $\frac{3\pi}4$.
\end{exercise}

\begin{exercise}[$\star\star\star$]\label{ex:v7c09d}
Theorem~\ref{thm:ball} says that the unit ball of $\R^3$ has volume $4\pi/3$. By Cavalieri, the slice of the unit ball of $\R^4$ at a height $t$ strictly between $-1$ and $1$ is a three-dimensional ball of radius $\sqrt{1-t^2}$. Compute $\lambda_4(B^4)$ from that description. The substitution $t=\sin\theta$ reduces the integral to a power of cosine. Check that the result is $\frac{\pi^2}2$, in agreement with Theorem~\ref{thm:ball}.
\end{exercise}

\section{Hints}

\begin{hint}{ex:fubinifail}
The inner integral in $y$ is the derivative of $y/(x^2+y^2)$. The integrand changes sign when the two variables are swapped. In polar coordinates the absolute value is $|\cos 2\theta|$ divided by $r^2$.
\end{hint}

\begin{hint}{ex:basel}
Integrate in $x$ first, with the substitution $x=u/\sqrt y$. Then integrate in $y$ first, by partial fractions in $y$. On $]1,\infty[$, the substitution $x=1/u$ folds the integral back to $]0,1[$. Expand $1/(1-x^2)$ as a geometric series and integrate term by term. The model integral of $x^{2k}$ times $-\ln x$ on the unit interval equals $1/(2k+1)^2$. Split $\zeta(2)$ into odd and even terms.
\end{hint}

\begin{hint}{ex:betagamma}
The inverse map sends $(u,v)$ to $(uv,\,u(1-v))$. It runs from $]0,\infty[\times\,]0,1[$ onto $]0,\infty[^2$, and the absolute value of its Jacobian is $u$.
\end{hint}

\begin{hint}{ex:gammasum}
The density of a sum of independent variables is the convolution of the densities. In the convolution integral over the segment from $0$ to $s$, substitute $x=sv$. The Beta integral you meet is Exercise~\ref{ex:betagamma}.
\end{hint}

\begin{hint}{ex:sphere}
For bounded Borel functions $g,h\ge 0$, compute the expectation of $g(|Z|)\,h(Z/|Z|)$ by Theorem~\ref{thm:polar}. For the ball, compute the probability that the radius is at most $r$ and the direction lies in $B$, and compare it with the normalised volume of the corresponding cone.
\end{hint}

\begin{hint}{ex:spherical3}
Expand the $3\times 3$ determinant along the last row. The map misses only a closed half-plane, which is null. For the integral, Theorem~\ref{thm:polar} is the short route.
\end{hint}

\begin{hint}{ex:shell}
In (a), the ball of radius $1-\varepsilon$ is a scaled copy of the unit ball. In (b), Cavalieri says that the slice at height $t$ is a ball of dimension $d-1$ and radius $\sqrt{1-t^2}$. Substitute $t=s/\sqrt d$. The bound $(1-s^2/d)^{(d-1)/2}\le e^{-s^2(d-1)/(2d)}$ dominates the integrand, and the ratio of successive ball volumes comes from Theorem~\ref{thm:ball}.
\end{hint}

\begin{hint}{ex:simplex}
The map $t\mapsto(s,x)$ is a diffeomorphism from $]0,\infty[^{n+1}$ onto $]0,\infty[\times\Delta_n$. Its inverse multiplies the simplex coordinates by the sum $s$, and puts the leftover mass on the last coordinate. Add the first $n$ rows of that inverse's Jacobian matrix to the last row. The determinant equals $s^n$.
\end{hint}

\begin{hint}{ex:frullani}
In (a), the ellipsoid is the image of the unit ball under a diagonal linear map. In (b), the integrand is positive, so Tonelli allows either order with no further check.
\end{hint}

\begin{hint}{ex:v7c09a}
At fixed $x$, the inner integral is the length $x^2$. At fixed $y$, the horizontal slice runs from $x=\sqrt y$ to $x=1$. The indicator is non-negative and the square has finite measure.
\end{hint}

\begin{hint}{ex:v7c09b}
The matrix of $T$ has columns $(2,0)$ and $(1,1)$, so the determinant equals $2$. At height $v$, the horizontal slice of the image runs from $u=v$ to $u=v+2$.
\end{hint}

\begin{hint}{ex:v7c09c}
On the open first quadrant the inverse is $r=\sqrt{x^2+y^2}$ and $\theta=\arctan(y/x)$. The columns of the Jacobian matrix are $(\cos\theta,\sin\theta)$ and $(-r\sin\theta,\,r\cos\theta)$.
\end{hint}

\begin{hint}{ex:v7c09d}
The slice at height $t$ has three-dimensional volume $\frac43\pi(1-t^2)^{3/2}$. After $t=\sin\theta$, the factor $\cos\theta$ from $dt$ produces $\cos^4\theta$ on $[-\pi/2,\pi/2]$. Rewrite $\cos^2\theta$ by the double-angle formula as $(1+\cos 2\theta)/2$, and square that identity.
\end{hint}

\section{Solutions}

\begin{solution}{ex:fubinifail}
\idea The integrand is a derivative in $y$, so the inner integral is explicit. Swapping the variables flips the sign, and the other order comes for free. Two different finite answers mean that $f$ cannot be integrable. Polar coordinates make the non-integrable spike at the origin visible.
\writeup Fix $x\in\,]0,1[$. The map $y\mapsto y/(x^2+y^2)$ has derivative
\[
\frac{(x^2+y^2)-2y^2}{(x^2+y^2)^2}=\frac{x^2-y^2}{(x^2+y^2)^2}=f(x,y).
\]
So $\int_0^1 f(x,y)\,dy=1/(1+x^2)$. A second integration gives $\int_0^1 dx/(1+x^2)=\arctan 1=\pi/4$. The identity $f(y,x)=-f(x,y)$ turns the other repeated integral into $-\pi/4$. Fubini's theorem would have forced the two numbers to agree, so $f$ is not integrable on $]0,1[^2$.

The square contains the quarter disc of radius $1$ in the open first quadrant. On that quarter disc the absolute value of $f$ equals $|\cos 2\theta|$ divided by $r^2$. The polar integral of the absolute value is
\[
\int_0^{\pi/2}\!\!\int_0^1\frac{|\cos 2\theta|}{r}\,dr\,d\theta,
\]
which factors as a positive angular integral times $\int_0^1 dr/r=\infty$. The positive part, near the $x$-axis, and the negative part, near the $y$-axis, each have infinite integral. Each order of integration subtracts those two infinities in a different way.
\end{solution}

\begin{solution}{ex:basel}
\idea The integrand is positive, so Tonelli allows both orders with no integrability check. One order produces $\pi^2/2$. The other produces a series. Every later swap is a sum of non-negative terms, so monotone convergence applies.
\writeup The integrand is continuous and positive on $]0,\infty[^2$. Tonelli says that both repeated integrals equal $I$, whether $I$ is finite or not.

Integrate in $x$ first. For $y>0$, the substitution $x=u/\sqrt y$ gives
\[
\int_0^\infty\frac{dx}{1+x^2 y}=\frac1{\sqrt y}\cdot\frac\pi2.
\]
Then
\[
I=\frac\pi2\int_0^\infty\frac{dy}{\sqrt y\,(1+y)}.
\]
The substitution $y=w^2$ turns the integral into $\int_0^\infty 2\,dw/(1+w^2)=\pi$, so $I=\pi^2/2$.

Integrate in $y$ first. For $x>0$ with $x\ne 1$, partial fractions give
\[
\frac1{(1+y)(1+x^2 y)}=\frac1{1-x^2}\Bigl(\frac1{1+y}-\frac{x^2}{1+x^2 y}\Bigr).
\]
The antiderivative is $\frac1{1-x^2}\ln\frac{1+y}{1+x^2 y}$. At the upper end the logarithm tends to $\ln(1/x^2)=-2\ln x$, and at the lower end it vanishes. So
\[
\int_0^\infty\frac{dy}{(1+y)(1+x^2 y)}=\frac{-2\ln x}{1-x^2}.
\]
The point $x=1$ is null, and therefore $\pi^2/2=\int_0^\infty\frac{-2\ln x}{1-x^2}\,dx$.

On $]1,\infty[$, substitute $x=1/u$. Then $dx=du/u^2$ after the limits are swapped, and $\ln x=-\ln u$. The integrand becomes
\[
\frac{-2(-\ln u)}{1-u^{-2}}\cdot\frac1{u^2}=\frac{-2\ln u}{1-u^2}.
\]
The two halves of the integral are equal, and each equals $\pi^2/4$. Dividing by $2$ gives
\[
\int_0^1\frac{-\ln x}{1-x^2}\,dx=\frac{\pi^2}8.
\]

On $]0,1[$, expand $1/(1-x^2)=\sum_{k\ge 0}x^{2k}$. The series $\sum_{k\ge 0}x^{2k}(-\ln x)$ has non-negative terms, so monotone convergence lets you integrate term by term. Integration by parts, with the boundary term $x^{2k+1}(-\ln x)/(2k+1)$ vanishing at both ends, gives
\[
\int_0^1 x^{2k}(-\ln x)\,dx=\int_0^1\frac{x^{2k}}{2k+1}\,dx=\frac1{(2k+1)^2}.
\]
Hence $\sum_{k\ge 0}1/(2k+1)^2=\pi^2/8$. Split $\zeta(2)$ into odd denominators and even denominators. The even terms are $\zeta(2)/4$, so $\zeta(2)=\pi^2/8+\zeta(2)/4$. Solving, $\zeta(2)=\pi^2/6$.
\end{solution}

\begin{solution}{ex:betagamma}
\idea In the product of the two Gamma integrals, the exponential $e^{-(s+t)}$ sees only the sum $s+t$. Taking that sum as a coordinate splits the integral into a Gamma factor in the sum and a Beta factor in the proportion $s/(s+t)$.
\writeup Tonelli gives
\[
\Gamma(a)\Gamma(b)=\int_{]0,\infty[^2}s^{a-1}t^{b-1}e^{-s-t}\,ds\,dt.
\]
The map $\psi(u,v)=(uv,\,u(1-v))$ is a $C^1$ bijection from $]0,\infty[\times\,]0,1[$ onto $]0,\infty[^2$. Its inverse $(s,t)\mapsto\bigl(s+t,\,s/(s+t)\bigr)$ is $C^1$ on that open quadrant. The Jacobian matrix of $\psi$ is
\[
\begin{pmatrix}v&u\\1-v&-u\end{pmatrix},
\]
with determinant $-uv-u(1-v)=-u$. Theorem~\ref{thm:cov} inserts the absolute value $u$, and Tonelli separates the product:
\begin{align*}
\Gamma(a)\Gamma(b)
&=\int_0^\infty\!\!\int_0^1(uv)^{a-1}\bigl(u(1-v)\bigr)^{b-1}e^{-u}\,u\,dv\,du\\
&=\int_0^\infty u^{a+b-1}e^{-u}\,du\cdot\int_0^1 v^{a-1}(1-v)^{b-1}\,dv\\
&=\Gamma(a+b)\,B(a,b).
\end{align*}
\end{solution}

\begin{solution}{ex:gammasum}
\idea The density of a sum of independent random variables is the convolution of their densities. For two Gamma densities the convolution integral is a Beta integral, and Exercise~\ref{ex:betagamma} evaluates it.
\writeup Write $f_a$ for the Gamma density of parameter $a$, equal to $t^{a-1}e^{-t}/\Gamma(a)$ on the positive half-line and equal to $0$ elsewhere. The sum of the two independent variables has density equal to the convolution $(f_a*f_b)(s)$, the integral of $f_a(x)f_b(s-x)$. For $s\le 0$ the integrand vanishes. For $s>0$ the substitution $x=sv$ gives
\[
(f_a*f_b)(s)=\frac{e^{-s}s^{a+b-1}}{\Gamma(a)\Gamma(b)}\int_0^1 v^{a-1}(1-v)^{b-1}\,dv,
\]
which equals $f_{a+b}(s)$ by Exercise~\ref{ex:betagamma}. An exponential random variable of parameter $1$ is Gamma of parameter $1$, because $\Gamma(1)=1$. By induction, the sum of $n$ independent copies has density $s^{n-1}e^{-s}/(n-1)!$ on $]0,\infty[$, since $\Gamma(n)=(n-1)!$. That is the waiting time for the $n$th event of a Poisson process of rate $1$ (Volume~4).
\end{solution}

\begin{solution}{ex:sphere}
\idea The Gaussian density depends on $z$ only through $|z|$. In polar coordinates it is a function of the radius times the constant function $1$ of the direction. The two factors separate, which is independence, and the direction is distributed exactly as the normalised surface measure. For the ball, the radius has to be reweighted so that its density is proportional to $r^{d-1}$.
\writeup Let $g$ and $h$ be bounded Borel functions, $g\ge 0$ on $]0,\infty[$ and $h\ge 0$ on $S=S^{d-1}$. Theorem~\ref{thm:polar} gives
\begin{align*}
\mathbb E\Bigl[g(|Z|)\,h\Bigl(\frac{Z}{|Z|}\Bigr)\Bigr]
&=\int_0^\infty\!\!\int_S g(r)\,h(\omega)\,\frac{e^{-r^2/2}}{(2\pi)^{d/2}}\,d\sigma(\omega)\,r^{d-1}\,dr\\
&=c\int_0^\infty g(r)\,e^{-r^2/2}r^{d-1}\,dr\cdot\int_S h\,\frac{d\sigma}{\sigma(S)},
\end{align*}
with $c=\sigma(S)/(2\pi)^{d/2}$. The constant $c$ is already inside the first factor. Taking $h=1$ makes the second factor equal to $1$, so the first factor equals $\mathbb E[g(|Z|)]$. Taking $g=1$ makes the first factor equal to $1$, because the Gaussian density integrates to $1$, and the left side is then $\mathbb E[h(Z/|Z|)]$. Therefore
\[
\mathbb E\bigl[g(|Z|)\,h(Z/|Z|)\bigr]=\mathbb E[g(|Z|)]\,\mathbb E[h(Z/|Z|)]
\]
whenever $g$ and $h$ are indicators, which is independence, and the law of $Z/|Z|$ is $\sigma/\sigma(S)$.

Now let $\omega=Z/|Z|$ and let $X=U^{1/d}\omega$, with $U$ independent of $\omega$. For $0<r\le 1$ and Borel $B\subseteq S$,
\[
\mathbb P\Bigl(|X|\le r,\ \frac{X}{|X|}\in B\Bigr)=\mathbb P(U\le r^d)\,\frac{\sigma(B)}{\sigma(S)}=r^d\frac{\sigma(B)}{\sigma(S)}.
\]
For the uniform law on $B^d$, Theorem~\ref{thm:polar} gives the same probability, because the cone out to radius $r$ has volume $\sigma(B)\,r^d/d$ and the ball has volume $\sigma(S)/d$. The sets $\{s\omega:0<s\le r,\ \omega\in B\}$ correspond under $\Phi$ to rectangles $]0,r]\times B$. Those rectangles form a $\pi$-system, the whole punctured ball corresponds to the rectangle $]0,1]\times S$, and both laws are finite. Theorem~\ref{thm:uniq6} says the two laws agree on the Borel sets of $B^d\setminus\{0\}$, hence everywhere.
\end{solution}

\begin{solution}{ex:spherical3}
\idea Compute the $3\times 3$ Jacobian directly, then read the radial integral off Theorem~\ref{thm:polar} with $\sigma(S^2)=4\pi$.
\writeup The map is $C^1$ on the open rectangle. Its inverse sends $x\ne 0$ to $r=|x|$, to $\theta=\arccos(x_3/|x|)$, and to the angle $\phi\in\,]0,2\pi[$ of the pair $(x_1,x_2)$. That inverse is $C^1$ off the closed half-plane $\{x_2=0,\ x_1\ge 0\}$. The half-plane is null, and it is exactly the set the image misses. The Jacobian matrix has rows
\[
\begin{gathered}
(\sin\theta\cos\phi,\ r\cos\theta\cos\phi,\ -r\sin\theta\sin\phi),\\
(\sin\theta\sin\phi,\ r\cos\theta\sin\phi,\ r\sin\theta\cos\phi),\\
(\cos\theta,\ -r\sin\theta,\ 0).
\end{gathered}
\]
Expand along the last row. The entry $\cos\theta$ has cofactor $r^2\cos\theta\sin\theta$, and the product of entry and cofactor is $r^2\cos^2\theta\sin\theta$. The entry $-r\sin\theta$ has cofactor $-r\sin^2\theta$, and the product of entry and cofactor is $r^2\sin^3\theta$. Adding these two contributions and factoring,
\[
J=r^2\sin\theta\,(\cos^2\theta+\sin^2\theta)=r^2\sin\theta.
\]
On the open rectangle the sine is positive, so $J>0$. The volume of the unit ball is
\[
\int_0^1 r^2\,dr\int_0^\pi\sin\theta\,d\theta\int_0^{2\pi}d\phi=\frac13\cdot 2\cdot 2\pi=\frac43\pi.
\]
The area of the sphere is the same integral without the factor $r^2$ and without the integral in $r$, which equals $4\pi$. That matches $\sigma(S^2)$. Finally Theorem~\ref{thm:polar} gives
\[
\int_{\R^3}\frac{e^{-|x|}}{|x|}\,dx=4\pi\int_0^\infty e^{-r}\,r\,dr=4\pi\,\Gamma(2)=4\pi.
\]
\end{solution}

\begin{solution}{ex:shell}
\idea Volume scales like the $d$-th power of the radius, so a slightly smaller ball is negligible once $d$ is large. For the slab, Cavalieri reduces the fraction to a one-dimensional integral of $(1-t^2)^{(d-1)/2}$. On the scale $t=s/\sqrt d$ that integrand looks like $e^{-s^2/2}$.
\writeup (a) The ball of radius $1-\varepsilon$ has volume $(1-\varepsilon)^d\lambda_d(B^d)$, by the scaling in Theorem~\ref{thm:lebesgue}. The shell between radii $1-\varepsilon$ and $1$ therefore has fraction $1-(1-\varepsilon)^d$, which tends to $1$.

(b) For $|t|<1$, the slice $\{x'\in\R^{d-1}:(t,x')\in B^d\}$ is the ball of radius $\sqrt{1-t^2}$. Write $\omega_d=\lambda_d(B^d)$. Cavalieri says that the fraction of volume in the slab is
\[
p_d=\frac{\omega_{d-1}}{\omega_d}\int_{-c/\sqrt d}^{c/\sqrt d}(1-t^2)^{(d-1)/2}\,dt,
\]
provided $d>c^2$, so that the slab sits inside the ball. The substitution $t=s/\sqrt d$ produces
\[
p_d=\frac{\omega_{d-1}}{\sqrt d\,\omega_d}\int_{-c}^c\Bigl(1-\frac{s^2}d\Bigr)^{(d-1)/2}ds.
\]
For each fixed $s$, the integrand tends to $e^{-s^2/2}$, because
\[
\frac{d-1}2\ln\Bigl(1-\frac{s^2}d\Bigr)=-\frac{s^2}2\cdot\frac{d-1}d+O\Bigl(\frac1d\Bigr).
\]
On the finite interval $[-c,c]$ the integrand is at most $1$ for all large $d$. Dominated convergence (Theorem~\ref{thm:dct6}) passes the limit inside the integral.

The same computation with $c$ replaced by $\sqrt d$ describes the whole ball, whose fraction is $1$:
\[
1=\frac{\omega_{d-1}}{\sqrt d\,\omega_d}\int_{-\sqrt d}^{\sqrt d}\Bigl(1-\frac{s^2}d\Bigr)^{(d-1)/2}ds.
\]
The inequality $1-u\le e^{-u}$ gives the bound $e^{-s^2(d-1)/(2d)}$. For $d\ge 2$ the exponent is at most $-s^2/4$, because $(d-1)/d\ge 1/2$. Extending the integrand by $0$ outside $[-\sqrt d,\sqrt d]$, dominated convergence with the integrable dominant $e^{-s^2/4}$ shows that the integral tends to $\int_\R e^{-s^2/2}\,ds=\sqrt{2\pi}$. Hence $\omega_{d-1}/(\sqrt d\,\omega_d)$ tends to $1/\sqrt{2\pi}$, and
\[
p_d\to\int_{-c}^c e^{-s^2/2}\,\frac{ds}{\sqrt{2\pi}}.
\]
Any hyperplane through the origin can replace $\{x_1=0\}$, by the rotation invariance in Theorem~\ref{thm:polar}. Most of the volume lies within distance $3/\sqrt d$ of every equator, since the standard normal distribution puts only a small mass outside $[-3,3]$.
\end{solution}

\begin{solution}{ex:simplex}
\idea As in Exercise~\ref{ex:betagamma}, the exponential depends on the variables only through their sum $s$. The proportions $x_i=t_i/s$ live in the simplex. The change of variables splits the product of Gamma integrals into a Gamma factor in $s$ and Dirichlet's integral over the simplex.
\writeup Define $\psi(s,x)$ on the half-line times the simplex by sending $(s,x)$ to the $(n+1)$-tuple whose first $n$ coordinates are $s x_i$ and whose last coordinate is $s$ times the leftover $1-\sum x_i$. This is a $C^1$ bijection onto the positive orthant. The inverse, which records the sum of the coordinates and then the first $n$ proportions, is $C^1$ on that orthant. Order the coordinates of the domain as $(x_1,\dots,x_n,s)$. For $i\le n$, row $i$ of the Jacobian matrix has an $s$ in column $i$ and an $x_i$ in the last column. The last row begins with $-s$ repeated $n$ times and ends with the leftover factor $1-\sum x_i$. Add the first $n$ rows to the last row. The last row becomes $(0,\dots,0,1)$, the diagonal entries are $s,\dots,s,1$, and the matrix is triangular. Row operations of this kind do not change the determinant, so $|J_\psi|=s^n$, using $s>0$.

Let $A$ be the sum of the parameters $a_1,\dots,a_{n+1}$. Write $p(x)$ for the integrand of Dirichlet's integral, the product of the powers $x_i^{a_i-1}$ for $i\le n$ and of the leftover factor $(1-\sum_i x_i)^{a_{n+1}-1}$. Along the point $t=\psi(s,x)$, the monomial $\prod_i t_i^{a_i-1}$ equals $s^{A-n-1}p(x)$. Theorem~\ref{thm:cov} and Tonelli give
\begin{align*}
\prod_i\Gamma(a_i)
&=\int_0^\infty\!\!\int_{\Delta_n}s^{A-n-1}p(x)\,e^{-s}\,s^n\,dx\,ds\\
&=\Gamma(A)\int_{\Delta_n}p(x)\,dx.
\end{align*}
When every parameter $a_i$ equals $1$, the integral of $p$ is the volume of the simplex, and the Gamma values collapse to $1/n!$. The same proportions $(t_1/s,\dots,t_{n+1}/s)$, for independent Gamma variables, follow the Dirichlet law on the simplex.
\end{solution}

\begin{solution}{ex:frullani}
\idea In (a) the set is a linear image of a set whose volume is already known. In (b) the integrand $(e^{-ax}-e^{-bx})/x$ is itself an integral of $e^{-xy}$ over $y\in[a,b]$, and a positive function may be integrated in either order.
\writeup (a) Let $T$ be the diagonal matrix with diagonal entries $a$, $b$ and $c$. A point $(x,y,z)$ lies in the ellipsoid if and only if $T^{-1}(x,y,z)=(x/a,\,y/b,\,z/c)$ lies in the unit ball $B^3$. The ellipsoid is $T(B^3)$. Theorems~\ref{thm:lineardet} and~\ref{thm:ball} give the volume $|\det T|\,\lambda_3(B^3)=\frac43\pi abc$.

(b) The map $(x,y)\mapsto e^{-xy}$ is continuous and positive on $]0,\infty[\times[a,b]$. Tonelli (Theorem~\ref{thm:tonelli}) says the two repeated integrals agree. Integrating in $x$ first, $\int_0^\infty e^{-xy}\,dx=1/y$ for each $y>0$, and $\int_a^b dy/y=\ln(b/a)$. Integrating in $y$ first, $\int_a^b e^{-xy}\,dy=(e^{-ax}-e^{-bx})/x$ for each $x>0$. Therefore
\[
\int_0^\infty\frac{e^{-ax}-e^{-bx}}{x}\,dx=\ln\frac ba.
\]
Each separate piece $\int_0^\infty e^{-ax}/x\,dx$ diverges, so the integral cannot be split into a difference of two convergent integrals. Tonelli treats the difference as one positive function.
\end{solution}

\begin{solution}{ex:v7c09a}
\idea The indicator is non-negative and the domain is a square of finite measure, so Tonelli applies before any calculation. One order sees vertical slices of length $x^2$. The other sees horizontal slices running from $\sqrt y$ to $1$.
\writeup The function $\mathbf 1_Q$ is Borel and non-negative, and both factors are Lebesgue measure on a finite interval, hence $\sigma$-finite. Theorem~\ref{thm:tonelli} says the two repeated integrals both equal $\lambda_2(Q)$.

At fixed $x\in\,]0,1[$, the slice in $y$ is the interval $]0,x^2[$, of length $x^2$. So the outer integral is $\int_0^1 x^2\,dx=1/3$. At fixed $y\in\,]0,1[$, one has $y<x^2$ precisely when $\sqrt y<x<1$, an interval of length $1-\sqrt y$. The outer integral is
\[
\int_0^1\bigl(1-\sqrt y\bigr)\,dy=\Bigl[y-\frac23 y^{3/2}\Bigr]_0^1=1-\frac23=\frac13.
\]
The two numbers agree, as Tonelli promised.
\end{solution}

\begin{solution}{ex:v7c09b}
\idea The determinant is the area factor, by Theorem~\ref{thm:lineardet}. The same area is the integral of the slice lengths, and every horizontal slice of this particular image has length $2$.
\writeup The matrix of $T$, with columns $(2,0)$ and $(1,1)$, has determinant $2$. The unit square has area $1$, so Theorem~\ref{thm:lineardet} multiplies that area by $2$.

For the slices, let $(u,v)=T(x,y)=(2x+y,\,y)$ with $(x,y)\in[0,1]^2$. Then $v=y$ runs through $[0,1]$, and $u=2x+v$ runs through $[v,\,v+2]$. Conversely, every such $(u,v)$ comes from $y=v$ and $x=(u-v)/2\in[0,1]$. So $T(S)$ is the set in the statement. At height $v$ the horizontal slice is an interval of length $2$. Cavalieri, in the form $\lambda_2=\lambda\otimes\lambda$, gives
\[
\lambda_2(T(S))=\int_0^1 2\,dv=2.
\]
The two computations agree, and both match the area marked on Figure~\ref{fig:rw-det}.
\end{solution}

\begin{solution}{ex:v7c09c}
\idea The polar map on this rectangle stays away from the origin and away from the positive $x$-axis, so the usual inverse in the open first quadrant is smooth. The Jacobian is the same $r$ computed for the whole plane, and integrating it over the rectangle is the area.
\writeup The map $\varphi$ is $C^1$ on $U$. If $\varphi(r,\theta)=\varphi(r',\theta')$ with both points in $U$, the Euclidean norms give $r=r'$, and then the angle in $]0,\pi/2[$ gives $\theta=\theta'$. So $\varphi$ is injective. A point $(x,y)$ of the open quarter annulus has $r=\sqrt{x^2+y^2}\in\,]1,2[$ and $\theta=\arctan(y/x)\in\,]0,\pi/2[$, and $\varphi(r,\theta)=(x,y)$. Both coordinate functions of this inverse are $C^1$ on the quarter annulus, because $r>1$ keeps the square root away from $0$ and $x>0$ keeps the arctangent away from its vertical asymptote. Thus $\varphi$ is a $C^1$ diffeomorphism onto that open set.

The Jacobian matrix has columns $(\cos\theta,\sin\theta)$ and $(-r\sin\theta,\,r\cos\theta)$. Its determinant is
\[
r\cos^2\theta+r\sin^2\theta=r.
\]
Theorem~\ref{thm:cov}, applied to the constant function $1$, gives the area
\[
\int_1^2\int_0^{\pi/2} r\,d\theta\,dr=\frac\pi2\int_1^2 r\,dr=\frac\pi2\cdot\frac32=\frac{3\pi}4.
\]
The boundary of the closed quarter annulus is a finite union of smooth curves, hence null, so the open set and the closed quarter annulus have the same area.
\end{solution}

\begin{solution}{ex:v7c09d}
\idea Cavalieri writes the four-dimensional volume as an integral of three-dimensional volumes. Each slice is a ball whose radius is $\sqrt{1-t^2}$, so the integrand is a power of $1-t^2$. The substitution $t=\sin\theta$ turns that power into a power of cosine, which the double-angle formula integrates.
\writeup For $|t|<1$, the slice of $B^4$ at height $t$ is a ball in $\R^3$ of radius $\rho=\sqrt{1-t^2}$. Scaling multiplies three-dimensional volume by $\rho^3$, so the slice has volume
\[
\frac43\pi\,(1-t^2)^{3/2}.
\]
Cavalieri therefore gives
\[
\lambda_4(B^4)=\int_{-1}^1\frac43\pi\,(1-t^2)^{3/2}\,dt.
\]
Substitute $t=\sin\theta$, so $dt=\cos\theta\,d\theta$ and $\theta$ runs from $-\pi/2$ to $\pi/2$. On that interval the cosine is non-negative, so the power $(1-\sin^2\theta)^{3/2}$ simplifies to $\cos^3\theta$. Multiplying by the extra factor $\cos\theta$ from the differential produces $\cos^4\theta$. The integral equals
\[
\frac43\pi\int_{-\pi/2}^{\pi/2}\cos^4\theta\,d\theta.
\]
The double-angle formula $\cos 2\theta=2\cos^2\theta-1$ rearranges to $\cos^2\theta=(1+\cos 2\theta)/2$. Squaring,
\[
\cos^4\theta=\frac{1+2\cos 2\theta+\cos^2 2\theta}4.
\]
The same formula at angle $2\theta$ gives $\cos^2 2\theta=(1+\cos 4\theta)/2$, and therefore
\[
\cos^4\theta=\frac38+\frac12\cos 2\theta+\frac18\cos 4\theta.
\]
Integrate term by term from $-\pi/2$ to $\pi/2$. The constant term contributes $(3/8)\cdot\pi=3\pi/8$. The integrals of $\cos 2\theta$ and of $\cos 4\theta$ over that symmetric interval are zero, because their antiderivatives $\frac12\sin 2\theta$ and $\frac14\sin 4\theta$ vanish at both endpoints. Hence
\[
\lambda_4(B^4)=\frac43\pi\cdot\frac{3\pi}8=\frac{\pi^2}2,
\]
which is the value Theorem~\ref{thm:ball} gives in dimension $4$.
\end{solution}

\section*{Sources and further reading}
\addcontentsline{toc}{section}{Sources and further reading}

\begin{itemize}
\item G.~Fubini, ``Sugli integrali multipli'', \emph{Rend. Accad. Lincei} 16 (1907), 608--614, is the note behind the integrable case. L.~Tonelli, ``Sull'integrazione per parti'', \emph{Rend. Accad. Lincei} 18 (1909), 246--253, is the note behind the non-negative case.
\item G.~B.~Folland, \emph{Real Analysis}, 2nd ed., Wiley, 1999, Sections~2.5--2.7, develops product measures, the behaviour of Lebesgue measure under linear maps, a proof of the change of variables by cubes in the style of Lemma~\ref{lem:cube}, and polar coordinates in the style of Theorem~\ref{thm:polar}.
\item W.~Rudin, \emph{Real and Complex Analysis}, 3rd ed., McGraw--Hill, 1987, Chapters~7--8, proves the change of variables by a different route, through the differentiation of measures.
\item D.~Ritelli, ``Another proof of $\zeta(2)=\frac{\pi^2}6$ using double integrals'', \emph{Amer. Math. Monthly} 120 (2013), 642--645, is the source of the double-integral evaluation in Exercise~\ref{ex:basel}.
\item K.~Ball, ``An elementary introduction to modern convex geometry'', in \emph{Flavors of Geometry}, MSRI Publications 31, Cambridge, 1997, explains the geometry of high-dimensional balls and cubes, including the concentration of volume used in Exercise~\ref{ex:shell}.
\item G.~E.~P.~Box and M.~E.~Muller, ``A note on the generation of random normal deviates'', \emph{Ann. Math. Statist.} 29 (1958), 610--611, turns the polar splitting of a planar Gaussian vector into a way of simulating it.
\item The opening stories are told in R.~Netz and W.~Noel, \emph{The Archimedes Codex}, Weidenfeld \& Nicolson, 2007, from Heiberg's reading of the palimpsest in 1906 through its later imaging, and in A.~Alexander, \emph{Infinitesimal: How a Dangerous Mathematical Theory Shaped the Modern World}, Farrar, Straus and Giroux, 2014, which follows the dispute over Cavalieri's indivisibles.
\end{itemize}
